\documentclass[11pt,a4paper]{article}
\usepackage[margin=25mm]{geometry}
\usepackage{amsmath,amssymb,graphicx,booktabs,hyperref,listings,xcolor}
\hypersetup{colorlinks=true,linkcolor=blue,urlcolor=blue}
\lstset{basicstyle=\ttfamily\scriptsize,breaklines=true,frame=single,columns=fullflexible}
\title{EN3563 Robotics\\GP7 Pick-and-Place Simulation}
\author{Team members and index numbers: \textbf{TO BE COMPLETED}}
\date{\today}
\begin{document}
\maketitle
\begin{abstract}
A CAD-based six-axis GP7 manipulator performs a pick-and-place task in ROS 2 Jazzy and Gazebo Harmonic. Numerical inverse kinematics generates a constant-orientation Cartesian trajectory with quintic timing. The corrected CAD gripper approaches, lifts and releases a cube on a destination platform. Grasp holding uses an explicitly idealized detachable fixed joint. Simulator measurements verify lift, orientation preservation and final placement. This draft distinguishes demonstrated simulation results from unresolved manufacturer calibration.
\end{abstract}
\section{Task and scene}
The task is to transfer a $50$ mm cube, Box B, while preserving its orientation. Box A measures $0.30\times0.30\times0.20$ m. The original assignment destination $(1,1,0)$ has horizontal distance $1.414$ m from the robot, exceeding the nominal GP7 reach $0.927$ m. The destination was relocated by user request. \textbf{Instructor approval must be obtained and recorded.}
\begin{table}[h]\centering
\begin{tabular}{ll}\toprule Item & World center position (m)\\\midrule
Box A & $(0.50,0.50,0.10)$\\
Box B source & $(0.40,-0.25,0.025)$\\
Box B destination & $(0.50,0.50,0.225)$\\
Robot spawn translation & $(0,0,0.002)$\\\bottomrule\end{tabular}
\caption{Current scene. Box A ground location is $(0.5,0.5,0)$.}
\end{table}
Box B has assumed mass $0.05$ kg and uniform-cube inertia. Its source was selected deterministically, not randomly sampled for each trial.
\section{Software and model preparation}
The host uses Ubuntu 24.04, ROS 2 Jazzy and Gazebo Harmonic 8.11. The supplied SolidWorks-to-URDF export contained CAD STL meshes and ROS1 package files. These were migrated to \texttt{ament\_cmake}, ROS2 Python launch and \texttt{ros\_gz\_sim}. No native SolidWorks conversion was performed on the Linux host.
The original arm meshes remain visual geometry. Arm collision geometry uses local mesh bounding boxes. The corrected export supplies CAD finger meshes with corrected geometry origins and opposite prismatic axes. A fixed world-to-base joint anchors the arm. The original export is preserved separately.
\section{Frames and forward kinematics}
\textbf{Insert a labeled coordinate-frame illustration before submission.} The six revolute joints connect base, shoulder, elbow and wrist links. Each exported joint has a fixed origin transform $O_i$ and a normalized local axis $\hat a_i$. Using column vectors,
\begin{equation}
T_{W,T}(q)=T_{W,B}\left[\prod_{i=1}^{6}O_i\,R(\hat a_i,q_i)\right]T_{6,L6}T_{L6,T}.
\end{equation}
URDF origin rotations use roll--pitch--yaw. The CAD-derived TCP translation in Link 6 is approximately
\begin{equation}
p_{6,T}=(0.0051601583,\;0.0014121731,\;0.1053659958)^\mathsf T\;\mathrm{m}.
\end{equation}
The corrected Link 6 to TCP transform is read from the active simulator URDF. The grasp TCP is 10 mm above cube center. Five non-unit axis vectors were normalized without altering direction. Manufacturer encoder zero offsets and signs have not been completely established.
\subsection{Standard D--H construction}
The exact exported axis lines admit the standard representation
\begin{equation}
T_{W,T}=B\prod_{i=1}^{6}\left[R_z(q_i+\theta_{0i})T_z(d_i)T_x(a_i)R_x(\alpha_i)\right]E.
\end{equation}
Table~\ref{tab:dh} reports approximate display values; the derivation script emits full precision. The second and third axes differ by approximately $1.44\times10^{-8}$ rad. Their common normal is extremely distant, producing large cancelling offsets. \textbf{These are frame artifacts, not physical arm lengths.} The stable URDF product is used for runtime FK.
\begin{table}[h]\centering\small
\begin{tabular}{rrrrr}\toprule$i$&$\theta_{0i}$ (rad)&$d_i$ (m)&$a_i$ (m)&$\alpha_i$ (rad)\\\midrule
1&1.57470927&-0.12674&2.4074124e-35&1.5707927\\
2&-3.14159083&30811506.1&0.00081185985&1.4386509e-08\\
3&-1.52558307&-30811506.1&-0.039999777&1.5707963\\
4&-3.14033776&-0.442001848&3.7220432e-07&1.5707963\\
5&-3.13788375&2.72995319e-08&1.1122742e-16&1.5707926\\
6&0.00000000&5.55111512e-17&-1.110223e-16&-3.5236571e-19\\
\\bottomrule
\end{tabular}\caption{Exact-axis D--H model, rounded only for display. See \texttt{DH.md} and the derivation output for full precision.}\label{tab:dh}\end{table}
The DH link rows are constant. The corrected DH-to-Link6 and Link6-to-TCP matrices are generated by \texttt{scripts/derive\_dh.py} from the active corrected-gripper URDF. Their composition forms the fixed end transform.
\section{Numerical inverse kinematics}
SciPy bounded nonlinear least squares minimizes the residual
\begin{equation}
r(q)=\begin{bmatrix}p(q)-p_d\\0.3\,\operatorname{rotvec}(R_dR(q)^\mathsf T)\end{bmatrix}.
\end{equation}
The factor $0.3$ is a numerical weighting, not a physical parameter. Multiple initial seeds and square-compatible fixed gripper yaws are searched. A branch is accepted only after checking the complete path, not just the pregrasp pose. Subsequent samples use the preceding solution as the initial guess. This is numerical IK; no closed-form analytical solution is claimed.
\section{Planning and control}
Each phase uses $161$ samples over $8$ s with quintic time scaling
\begin{equation}s(u)=10u^3-15u^4+6u^5,\qquad u=t/8.
\end{equation}
The initial pregrasp phase interpolates in joint space. Grasp, lift, transfer, place and retreat use Cartesian position interpolation with fixed orientation. Joint velocities are estimated numerically and set to zero at phase endpoints. Joint limits, sample-to-sample continuity, velocity ceilings and sampled collisions are checked. These checks do not certify continuous swept volume or self collision.
\begin{table}[h]\centering\begin{tabular}{ll}\toprule Phase & Purpose\\\midrule
Pregrasp & Move above source\\Grasp & Descend without disturbing cube\\Lift & Raise cube approximately $0.30$ m\\Transfer & Move above destination\\Place & Lower to destination\\Retreat & Withdraw after release\\\bottomrule\end{tabular}\caption{Motion phases.}\end{table}
\texttt{gz\_ros2\_control} connects ideal position interfaces to arm and gripper trajectory controllers. \texttt{FollowJointTrajectory} actions execute the plan; the joint-state broadcaster supplies measured joint states. This is not a calibrated torque-control system.
\section{Gripper clearance and attachment}
The corrected CAD palm originally intersected the cube during descent. The failed run rotated the cube by about $0.48$ rad. Raising the grasp TCP from $2$ mm to $10$ mm above cube center removed the interference and gave approximately $3.4$ mm conservative palm clearance. The same offset is used at placement, preserving the cube destination. Finger closing displacement is $0.023$ m per finger.
A disturbance guard aborts before attachment if approach moves the cube by more than $3$ mm or rotates it by more than $0.02$ rad. The installed detachable-joint plugin begins attached; the executor first detaches and requires acknowledgement before motion. After approach and closure, it attaches Box B to Link 6. Release removes the constraint and lets gravity settle the cube. No object teleport commands are used. \textbf{Holding is idealized adhesion, not verified frictional force closure.}
\section{Results}
The latest trial is recorded in \texttt{results/fixed\_grasp\_clearance\_run.json}. It passed the approach guard, lift, carry-orientation and released-placement checks.
\begin{table}[h]\centering\begin{tabular}{ll}\toprule Measurement & Result\\\midrule
Final cube center (m)&$(0.49999916,0.49998903,0.22499936)$\\
Position error&$1.1025\times10^{-5}$ m\\
Final orientation error&$3.8398\times10^{-7}$ rad\\
Maximum sampled carry orientation error&$1.9270\times10^{-6}$ rad\\
Lift&approximately $0.30$ m\\\bottomrule\end{tabular}\caption{Measured simulator outcomes, not hardware accuracy.}\end{table}
The small numerical errors reflect ideal actuation and attachment. Acceptance tolerances are $25$ mm final position and $0.15$ rad orientation, with at least $50$ mm observed lift. \textbf{Insert screenshots and a new video showing the corrected CAD gripper and current scene. Older recordings show earlier geometry or destinations; this latest run JSON validates corrected-gripper motion, but a new current-scene video is still needed.}
\section{Limitations and calibration}
The original exported joint ranges $[-3.14,3.14]$ rad and $2$ rad/s speed ceilings are not verified manufacturer-coordinate specifications. Published GP7 ranges cannot be applied without sign/zero mapping. Arm effort $200$ Nm is a simulation placeholder. Original mass totals approximately $20.30$ kg, compared with published totals $34$--$37$ kg. Inertias pass positive-definiteness and triangle checks, but physical material assignments and missing internal parts remain unresolved. No arbitrary mass scaling was applied.
\section{Reproduction}
From the Projects directory:
\begin{lstlisting}[language=bash]
source /opt/ros/jazzy/setup.bash
colcon build --base-paths robo2 --build-base .ros2_build \
 --install-base .ros2_install --symlink-install \
 --cmake-args -DPython3_EXECUTABLE=/usr/bin/python3
source .ros2_install/setup.bash
ros2 launch robo2 gazebo.launch.py
# Another sourced terminal, fresh scene for each run:
ros2 run robo2 pick_place.py --results robo2/results/my_run.json
\end{lstlisting}
\appendix
\section{Code excerpts}
The submitted source files are authoritative. These excerpts are read directly from them during compilation.
\subsection{Forward and inverse kinematics}
\lstinputlisting[language=Python,firstline=15,lastline=56]{../scripts/gp7_kinematics.py}
\subsection{Trajectory generation and checks}
\textbf{Review line ranges after future code edits.} See the complete generator in the submitted package for collision checking and branch search.
\subsection{Execution and velocity control}
\lstinputlisting[language=Python,firstline=80,lastline=123]{../scripts/pick_place.py}
\begin{thebibliography}{9}
\bibitem{task} EN3563 Robotics, Mini Project: Implementation of a Robot Arm, supplied task sheet.
\bibitem{manual} Yaskawa, MOTOMAN-GP8/AR700, GP7/AR900 Instructions, HW1483944, revision 11; supplied manual.
\bibitem{data} Yaskawa America, GP7 and GP8 Compact High-Speed Robots, DS-866, January 2021.
\bibitem{ros} ROS 2 Jazzy documentation, \url{https://docs.ros.org/en/jazzy/}.
\bibitem{gz} Gazebo Harmonic documentation, \url{https://gazebosim.org/docs/harmonic/}.
\end{thebibliography}
\end{document}
